% Part: sets-functions-relations
% Chapter: functions
% Section: basics

\documentclass[../../../include/open-logic-section]{subfiles}

\begin{document}

\olfileid{sfr}{fun}{bas}
\olsection{ప్రాథమిక విషయాలు}

\begin{explain}
ఇచ్చిన సమితిలోని ప్రతి \tecase{acc}{!!{element}} మరో ఇచ్చిన సమితిలోని
ఒక నిర్దిష్ట \tecase{dat}{!!{element}} పంపే ప్రతిచిత్రణాన్ని
\emph{ప్రమేయం} అంటాం. ఉదాహరణకు, $1$ను కలిపే ప్రక్రియ ఒక ప్రమేయాన్ని
నిర్వచిస్తుంది: ప్రతి సంఖ్య $n$ ఒకే ఒక్క సంఖ్య $n+1$కు ప్రతిచిత్రించబడుతుంది.

మరింత సాధారణంగా, ప్రమేయాలు జతలను, త్రికాలను మొదలైనవాటిని ఇన్‌పుట్లుగా
తీసుకుని ఏదో ఒక ఫలితాన్ని ఇస్తాయి. ప్రాథమిక అంకగణితం నుంచి మనకు ఎన్నో
ప్రమేయాలు పరిచితమే. ఉదాహరణకు, కూడిక, గుణకారం ప్రమేయాలు.
అవి రెండు సంఖ్యలను తీసుకుని మూడో సంఖ్యను ఇస్తాయి.

ఈ అమూర్త గణిత అర్థంలో ప్రమేయం ఒక \emph{నల్ల పెట్టె}:
ఏ ఇన్‌పుట్‌కు ఏ ఫలితం జతచేయబడుతుందన్నదే ముఖ్యం;
ఆ ఫలితాన్ని గణించే పద్ధతి కాదు.
\end{explain}

\begin{defn}[ప్రమేయం]
$A$లోని ప్రతి \tecase{acc}{!!{element}}, $B$లోని ఒక \tecase{dat}{!!{element}}
ప్రతిచిత్రించే ప్రతిచిత్రణమే \emph{ప్రమేయం} $f \colon A \to B$.

$A$ను $f$ యొక్క \emph{ప్రవేశం} అనీ, $B$ను $f$ యొక్క
\emph{సహప్రవేశం} అనీ అంటాం. $A$లోని !!{element}sను $f$ యొక్క
ఇన్‌పుట్లు లేదా \emph{ఆర్గ్యుమెంట్లు} అంటాం. $f$ ఒక ఆర్గ్యుమెంట్ $x$కు
జతచేసే $B$లోని \tecase{acc}{!!{element}}, ఆ ఆర్గ్యుమెంట్ $x$ వద్ద
\emph{$f$ యొక్క విలువ} అని పిలిచి $f(x)$ అని రాస్తాం.

$f$ యొక్క \emph{వ్యాప్తి} $\ran{f}$ అంటే, ఏదో ఒక ఆర్గ్యుమెంట్ వద్ద
$f$ విలువలుగా వచ్చే మూలకాలన్నీ గల సహప్రవేశపు ఉపసమితి;
$\ran{f} = \Setabs{f(x)}{x \in A}$.
\end{defn}

ప్రమేయాలను అర్థం చేసుకోవడానికి \olref{fig:function}లోని పటం ఉపయోగపడవచ్చు.
ఎడమవైపు దీర్ఘవృత్తం ప్రమేయపు \emph{ప్రవేశాన్ని}, కుడివైపు దీర్ఘవృత్తం
దాని \emph{సహప్రవేశాన్ని} సూచిస్తాయి. ప్రవేశంలోని ఒక \emph{ఆర్గ్యుమెంట్}
నుంచి సహప్రవేశంలోని దానికి సంబంధించిన \emph{విలువ} వైపు బాణం చూపుతుంది.

\begin{figure}
  \olasset{assets/diagrams/function.tikz}
  \caption{ఒక సమితిలోని ప్రతి \tecase{acc}{!!{element}} మరో సమితిలోని
    \tecase{dat}{!!a{element}} ప్రతిచిత్రించడమే ప్రమేయం.
    ప్రవేశంలోని ఆర్గ్యుమెంట్ నుంచి సహప్రవేశంలోని సంబంధిత విలువ వైపు బాణం చూపుతుంది.}
  \ollabel{fig:function}
\end{figure}

\begin{ex}
గుణకారం సహజ సంఖ్యల జతలను ఇన్‌పుట్లుగా తీసుకుని, సహజ సంఖ్యలను
ఫలితాలుగా ఇస్తుంది. కాబట్టి అది $\Nat \times \Nat$ (ప్రవేశం)
నుంచి $\Nat$ (సహప్రవేశం)కు వెళుతుంది. ప్రతి $n \in \Nat$ కూడా
$n \times 1$ కనుక, వ్యాప్తి కూడా $\Nat$ అవుతుంది.
\end{ex}

\begin{ex}
ప్రతి ఇన్‌పుట్‌ను---అంటే సహజ సంఖ్యల ప్రతి జతను---ఒకే ఫలితంతో
జతచేస్తుంది కనుక గుణకారం ఒక ప్రమేయం:
$\times \colon \Nat^2 \to \Nat$.
కానీ సహజ సంఖ్య $n$ను $x^2 = n$ అయ్యే వాస్తవ సంఖ్య $x$కు
ప్రతిచిత్రించడం ప్రమేయం కాదు. ఎందుకంటే ప్రతి ధన పూర్ణ సంఖ్య $n$కు
రెండు వర్గమూలాలు ఉంటాయి: $\sqrt{n}$, $-\sqrt{n}$.
ధన వర్గమూలాన్ని మాత్రమే ఫలితంగా ఇచ్చి దీన్ని ప్రమేయంగా చేయవచ్చు:
$\sqrt{\phantom{X}} \colon \Nat \to \Real$.
\end{ex}

\begin{ex}
ఒక తరగతిలోని ప్రతి విద్యార్థిని ఆ విద్యార్థి తుది గ్రేడుతో జతచేసే
సంబంధం ఒక ప్రమేయం: ఒకే తరగతిలో ఏ విద్యార్థికీ రెండు వేర్వేరు తుది
గ్రేడులు రావు. ప్రతి విద్యార్థిని వారి తల్లిదండ్రులతో జతచేసే సంబంధం
ప్రమేయం కాదు: విద్యార్థులకు తల్లిదండ్రులు ఎవరూ లేకపోవచ్చు, ఇద్దరు
ఉండవచ్చు, లేదా అంతకంటే ఎక్కువమంది ఉండవచ్చు.
\end{ex}

\begin{explain}
సాధ్యమైన ప్రతి ఆర్గ్యుమెంట్ వద్ద ప్రమేయపు విలువ ఏమిటో కచ్చితంగా
పేర్కొని ప్రమేయాలను నిర్వచించవచ్చు. ఒక సూత్రాన్ని ఇవ్వడం, విలువను
గణించే పద్ధతిని వివరించడం, లేదా ప్రతి ఆర్గ్యుమెంట్‌కు విలువను
జాబితాగా ఇవ్వడం ఇందుకు వేర్వేరు మార్గాలు. ఏ విధంగా నిర్వచించినా,
ప్రతి ఆర్గ్యుమెంట్‌కు ఒకే ఒక్క విలువను పేర్కొనేలా చూసుకోవాలి.
\end{explain}

\begin{ex}
$f(x) = x+1$ అయ్యేలా $f \colon \Nat \to \Nat$ను నిర్వచిద్దాం.
ఈ నిర్వచనం $f$ సహజ సంఖ్యలను తీసుకుని సహజ సంఖ్యలను ఇచ్చే ప్రమేయమని
పేర్కొంటుంది. ఒక సహజ సంఖ్య $x$ను ఇస్తే, $f$ దాని ఉత్తరవర్తి $x+1$ను
ఇస్తుందని చెబుతుంది. ఈ సందర్భంలో సహప్రవేశం $\Nat$, $f$ వ్యాప్తి కాదు:
సహజ సంఖ్య $0$ ఏ సహజ సంఖ్యకూ ఉత్తరవర్తి కాదు.
$f$ వ్యాప్తి అన్ని ధన పూర్ణ సంఖ్యల సమితి $\Int^{+}$.
\end{ex}

\begin{ex}\ollabel{examplefunext}
$g(x) = x+2-1$ అయ్యేలా $g \colon \Nat \to \Nat$ను నిర్వచిద్దాం.
దీని ప్రకారం $g$ సహజ సంఖ్యలను తీసుకుని సహజ సంఖ్యలను ఇచ్చే ప్రమేయం.
ఒక సహజ సంఖ్య $n$ ఇచ్చినప్పుడు, $g$, $x$ యొక్క ఉత్తరవర్తి యొక్క
ఉత్తరవర్తి యొక్క పూర్వవర్తిని, అంటే $x+1$ను, ఫలితంగా ఇస్తుంది.
\end{ex}

\begin{explain}
ఇప్పుడే వేర్వేరు \emph{నిర్వచనాలు} గల $f$, $g$ అనే రెండు ప్రమేయాలను
పరిశీలించాం. అయితే ఇవి \emph{ఒకే ప్రమేయం}. ఏ సహజ సంఖ్య $n$కైనా
$f(n) = n+1 = n+2-1 = g(n)$ అవుతుంది కదా.
మరో విధంగా చెప్పాలంటే, $f$, $g$లకు ఇచ్చిన నిర్వచనాలు వేర్వేరు
సమీకరణాల ద్వారా ఒకే ప్రతిచిత్రణాన్ని పేర్కొంటాయి. అంటే ప్రమేయాలకు
వర్తించే మూలకాధారిత సమానత్వ సూత్రంపై మనం అంతర్లీనంగా ఆధారపడుతున్నాం:
\[
  \text{ఒకవేళ }\forall x\, f(x) = g(x)\text{ అయితే }f = g
\]
ఇందుకు $f$, $g$ల ప్రవేశమూ సహప్రవేశమూ ఒకటే కావాలి.
\end{explain}

\begin{ex}
సందర్భాలవారీగా కూడా ప్రమేయాలను నిర్వచించవచ్చు.
ఉదాహరణకు, $h \colon \Nat \to \Nat$ను ఇలా నిర్వచించవచ్చు:
\[
h(x) =
\begin{cases}
  \frac{x}{2} & \text{$x$ సరి సంఖ్య అయితే} \\
  \frac{x+1}{2} & \text{$x$ బేసి సంఖ్య అయితే.}
\end{cases}
\]
ప్రతి సహజ సంఖ్య సరి లేదా బేసి కనుక, ఈ ప్రమేయపు ఫలితం ఎల్లప్పుడూ
సహజ సంఖ్యే అవుతుంది. సందర్భాలవారీగా ప్రమేయాన్ని నిర్వచిస్తే,
సాధ్యమైన ప్రతి ఇన్‌పుట్ కచ్చితంగా ఒకే సందర్భంలోకి రావాలని గుర్తుంచుకోండి.
కొన్నిసార్లు సందర్భాలు అన్ని అవకాశాలనూ ఆవరించి, పరస్పరం అతివ్యాప్తి
చెందవని నిరూపించాల్సి ఉంటుంది.
\end{ex}

\end{document}
